%!TEX program = uplatex
\RequirePackage{plautopatch}
\documentclass[uplatex,dvipdfmx]{jlreq}
\usepackage{amsmath,amssymb}

\pagestyle{empty}

\begin{document}

%======================================================================
% 表紙
%======================================================================
\begin{center}
  \vspace*{20pt}
  {\Huge\bfseries 岩澤理論と Fermat 予想}

  \vspace{1pc}
  {\Large 概要}
\end{center}

\vfill

\newpage

%======================================================================
% 講演１：加藤和也「Fermat の最終定理と岩澤理論」
%======================================================================
\noindent
{\large\bfseries Fermat の最終定理と岩澤理論}

\bigskip
\hspace*{\fill}{\large\bfseries 加藤　和也（東工大・理）}

\bigskip

Wiles によるフェルマーの最終定理の証明の核心や背景、岩澤理論との関係について、
できるだけわかりやすく説明したい。
\begin{enumerate}
  \item Wiles と岩澤理論と楕円曲線
  \item 楕円曲線と保型形式
  \item 保型形式と岩澤理論と Wiles
\end{enumerate}
の3つに分けて述べさせていただく。

\medskip
\noindent\textbf{\S1.}\quad Wiles と岩澤理論と楕円曲線

Wiles のこれまでの主な仕事を年代順にあげると次のようになる。
\begin{itemize}
  \item[(ア)] (1970年代)\quad 楕円曲線の有理点に関する Birch と Swinnerton-Dyer 予想についての
              画期的な仕事。（岩澤理論の方法による。）
  \item[(イ)] (1980年代)\quad 岩澤 Main Conjecture の証明と、その総実代数体への一般化。
  \item[(ウ)] フェルマーの最終定理に関わる今回の仕事。（岩澤理論の拡張による。）
\end{itemize}
こうしてみると、Wiles は岩澤理論を深化・拡張してきた人であった。
岩澤理論は、「イデアル類群と、ゼータ関数の $p$ 進世界でのふるまいの間には、
ふしぎな関係がある」というものである。
上の（ア）や（ウ）での Wiles の方法は、
「問題を $p$ 進世界へ持っていって考え、岩澤理論の方法を適用する」というものである。
Wiles の（ア）や（ウ）の仕事は楕円曲線に関わるものであるが、
Wiles 同様 Fermat も、楕円曲線の有理点や整数点を考察するのが好きな人であった。
Fermat や Wiles をひきつけた「楕円曲線の数論」というものについて、
Wiles の仕事（ア）をふりかえりつつ述べる。

\medskip
\noindent\textbf{\S2.}\quad 楕円曲線と保型形式

「楕円曲線と保型形式がふしぎにも対応する」という、谷山--志村（--Weil）予想を紹介し、
フェルマーの最終定理がこの予想に帰着されたこと、
この予想が「類体論の非可換版」のひとつと見なせること、を述べる。
Wiles はこの予想をほぼ証明することにより、フェルマーの最終定理を証明したのであった。

\medskip
\noindent\textbf{\S3.}\quad 保型形式と岩澤理論と Wiles

この予想をほぼ証明した Wiles の方法を紹介する。
これは、問題を $p$ 進世界に持っていって考え、
岩澤理論を拡張して考えることによるものであった。

\newpage

%======================================================================
% 講演２：百瀬文之「Langlands 予想への試み」
%======================================================================
\noindent
{\large\bfseries Langlands 予想への試み}

\bigskip
\hspace*{\fill}{\large\bfseries 百瀬　文之（中大・理工）}

\bigskip

数論の世界では、ゼータ関数なる方々が重きを成しており、
この分野の心を殆ど表現していると思われております。
このゼータ様について、ある意味で一応分かっているのは1次の場合だけです。
それは、代数体 $k$ のガロア群のアーベルな表現とか、
Hecke の表現に対応する場合です。
では、$\mathrm{GL}_n$（$n \geq 2$）への非アーベルに対応する場合はどうかと申しますと、
$\mathrm{GL}_n(\mathbf{C})$ への（有限位数の）表現の分りやすい場合は、
アーベル多様体 $A$ の素数冪分点への Galois 表現に対応するゼータ関数です。
それも $A$ が非常に特殊な形、いわゆる CM type の場合のみです。

一般の場合で、一応分っていた例は、保型関数に対応した場合です。
古典的な楕円型モジュラー保型形式には $\mathrm{GL}_2$ type の表現が対応します。
この場合について有名な谷山--志村予想（条件付き）を解いたのが Wiles さんたちです。
そして、それは即、Fermat 予想の証明を与えてくれました。
（この部分は、Frey と Ribet によります。）
Wiles さんの御弟子さんの Diamond 氏が条件をより弱い形にしました。
これに Ribet--Pyle の結果を用いると、様々な予想が（条件付きですが）解けます。
いわゆる、Modular Conjecture for $\mathrm{GL}_2$ type $/\mathbf{Q}$ です。
中には Ribet 氏の $\mathbf{Q}$-curves に対する予想、
Hashimoto 氏の QM-curves に対する予想が入っており、
これがかなりの予想で分ります。

藤原氏が総実体上の $\mathrm{GL}_2$ type について Wiles さんの結果を一般化しました
（もちろん条件付きですが）。
これからは $n \geq 3$ の場合を扱うべきです。
分りやすい場合として重さ2の保型形式に対応すると思われる
アーベル多様体のゼータ関数の分類をしてみました。
出来たらこの辺も話したく思います。

\newpage

%======================================================================
% 講演３：藤原一宏「楕円曲線と進化するシステム --- Euler 系から Taylor-Wiles 系へ」
%======================================================================
\noindent
{\large\bfseries 楕円曲線と進化するシステム --- Euler 系から Taylor-Wiles 系へ}

\bigskip
\hspace*{\fill}{\large\bfseries 藤原　一宏（名大・多元数理）}

\bigskip

1989年に Kolyvagin が "Euler system" という概念を生み出したとき整数論は大きく変わりました。
彼自身が構成したのは Heegner point の Euler 系で、Birch--Swinnerton-Dyer 予想の特別な場合、
つまり modular elliptic curve $E$ に対し
\[
  \mathrm{ord}_{s=1} L(E, s) \le 1
  \;\Rightarrow\;
  \mathrm{rank}\, E(\mathbf{Q}) = \mathrm{ord}_{s=1} L(E, s)
\]
を示しました。
（円分体論における岩澤の主予想（の別証）もこの手法で示せます。
後になって判ったことですが、実は Kummer の仕事の中に
Euler system のアイディアが埋もれていたのです。）

その後いくつかの Euler 系が生まれ、整数論において最も基本的な
「セルマー群（イデアル類群の一般化）の位数 $=$ $L$ 関数の値」
という関係を導いてきました。（Bloch--加藤予想と呼ばれる。）
Wiles はこの等式がガロア表現の研究に使えることに気づき、
Euler 系を作ることで谷山--志村予想を解決できると信じたのですが、
彼の最初の証明はまさにこの部分にギャップがあったのです。

その後の1994年10月の Taylor との共著論文のなかで、
Wiles は Euler 系を使うことなく
（$\mathbf{Q}$ 上の準安定楕円曲線に対する）谷山--志村予想の証明を完成しました。
といっても、Euler 系の思想は受け継がれています。
この Taylor と Wiles の仕事の中から整数論における新しいシステムが誕生したのです。
この新しいシステム --- Taylor-Wiles 系は作りやすく、幾何学的意味が非常に明快であり、
他分野の方でも理解しやすいものであると信じます。
もちろん、かつてなかったほど強力な結果も次々に生み出します。

この講演では名付け親として Taylor-Wiles 系の「力」の一端をお見せできればよいなあと
思っています。

\end{document}